\section{Background and Related Work}
\label{sec:background}

Two lines of work assign functions to attention heads. The first \emph{describes} a
head by its attention pattern and weights. The second \emph{intervenes} on the head
and measures what changes in the output. Semantic induction heads were defined in the
first way; we audit them with the tools of the second.

\subsection{Induction heads and semantic induction heads}
\label{sec:sih}

\paragraph{Induction heads.}
An attention head can be split into a QK circuit, which chooses where the head
attends, and an OV circuit, which determines what it writes
\citep{elhage2021mathematical}. An \emph{induction head} uses the QK circuit to attend
from a repeated token $[A]$ to the token $[B]$ that followed its earlier occurrence,
and the OV circuit to raise the logit of $[B]$. The behaviour is defined on controlled
repeated sequences, and its importance for in-context learning was established by
ablation \citep{olsson2022induction}. Appendix~\ref{app:ih} gives the fuller picture,
including retrieval heads \citep{wu2024retrieval}.

\paragraph{Semantic induction heads.}
\citet{ren2024semantic} keep the same QK/OV template but let the written token differ
from the attended one. Given a relation triplet $(t_s, r, t_o)$ in the context---a
\emph{head} token $t_s$, a relation $r$ and a \emph{tail} token $t_o$---a
\emph{semantic induction head} (SIH) is a head that, at a later current token $t_j$,
attends to $t_s$ and raises the logit of $t_o$ (Figure~\ref{fig:ri}). The
\emph{relation index} (RI) turns this into a score. Let $p^{h,j}=\mathrm{softmax}(x_j
W^h_{OV} W_U)$ be the vocabulary distribution obtained by pushing the current token's
input embedding $x_j$ through head $h$'s OV weights and the unembedding, restricted to
the tokens seen so far and centred at zero. Then
\begin{equation}
\label{eq:ri}
\mathrm{RI}_h \;=\; \operatorname*{mean}_{(T,j)\,:\,\text{$h$ attends to } t_s}
\;\frac{q^{h,j}_{t_o}}{\sum_{k\le j} q^{h,j}_{t_k}},
\end{equation}
where $q^{h,j}_t=\max(0,\,p^{h,j}_t-\bar p^{h,j})$ and the mean runs over all
(triplet, position) pairs at which the head's attention argmax is $t_s$ by a fixed
margin. The index is the share of the head's positive output mass that lands on the
tail token; heads with a high index for a relation are called SIHs. The authors also
report that in a model they train, the index of some heads rises at the training
stage where few-shot pattern discovery emerges, and conclude that SIHs play a crucial
role in in-context learning. Appendix~\ref{app:ri} restates the definition exactly.

\begin{figure}[t]
\centering
\begin{tikzpicture}[
  tok/.style={draw, rounded corners=1pt, minimum height=1.5em, inner sep=2.5pt, font=\small},
  lab/.style={font=\scriptsize, text=black!65},
  >=stealth, node distance=1.5pt]
\node[tok] (t1) {Anna};
\node[tok, right=of t1] (t2) {is};
\node[tok, right=of t2] (t3) {the};
\node[tok, right=of t3] (t4) {mother};
\node[tok, right=of t4] (t5) {of};
\node[tok, right=of t5] (t6) {Ben};
\node[tok, right=of t6] (t7) {.};
\node[tok, right=of t7] (t8) {$\cdots$};
\node[tok, right=of t8, fill=black!8] (t9) {$t_j$};
\node[lab, below=1pt of t1] {$t_s$ (head)};
\node[lab, below=1pt of t6, xshift=-11pt] {$t_o$ (tail)};
\node[lab, below=1pt of t9] {current};
% QK: arc above the tokens
\draw[->, thick] (t9.north) to[out=120, in=60, looseness=0.45]
  node[lab, above, pos=.5, yshift=1pt] {QK circuit: attention argmax on $t_s$} (t1.north);
% OV: down from t_j, project, raise the tail token
\node[tok, dashed, below=1.05cm of t5, inner sep=3pt] (voc)
  {OV circuit: $\mathrm{softmax}(x_j W^h_{OV} W_U)$ over context tokens};
\draw[->, thick] (t9.south) -- ++(0,-0.62) -| (voc.east);
\draw[->, thick, dotted] (voc.north -| t6) -- node[lab, right=3pt, pos=.35] {raises $t_o$?} (t6.south |- t6.south) ;
\end{tikzpicture}
\caption{The relation index. For a relation (Anna, mother-of, Ben) stated in the
context, a head scores on a later current token $t_j$ if its attention argmax falls
on the head token $t_s$ and the vocabulary projection of its OV output, computed from
the \emph{input embedding} of $t_j$, puts a large share of its mass on the tail token
$t_o$. The projection is read at $t_j$, not at the attended position.}
\label{fig:ri}
\end{figure}

\paragraph{What the original study does not establish.}
The SIH evidence is descriptive. \emph{(i)} No head is intervened on, so it is not
tested whether high-RI heads are needed for the model to use a relation, or whether
the heads that are needed have a high RI. \emph{(ii)} The training-dynamics claim
rests on co-timing between a selected set of rising indices and a behavioural
transition, in a separately trained model; it is not shown that the heads whose
index rises are the heads that carry the ability, nor that they are the final SIHs.
\emph{(iii)} The index projects the \emph{raw embedding} of the current token, an
approximation borrowed from attention-only models. In a deep model the head's value
input is a contextual residual, and what the head writes at $t_j$ depends on the
positions it attends to, not on the embedding at $t_j$. There is no evident reason
for such a score to track the head's contribution. \emph{(iv)} The index has no
control (the tail competes only with all earlier tokens) and no null distribution, so
``high'' has no reference. Appendix~\ref{app:ri} details each point against the
original text. We address (i) by causal measurement of every head
(\S\ref{sec:method-causal}); (iii) and (iv) by re-implementing the RI with matched
controls and a null calibration (\S\ref{sec:method-ri}) and by testing whether the
head's actual contextual input repairs the score (\S\ref{sec:results-char}). We do not
study training dynamics; (ii) remains a limitation of the original evidence.
